\documentclass[10pt,a4paper]{amsart}
\usepackage[margin=19mm]{geometry}
\usepackage{amsmath,amssymb,mathtools}
\usepackage[scr=rsfs,cal=euler]{mathalfa}
\usepackage{xfrac}
\usepackage[unicode,hidelinks,bookmarks=false]{hyperref}
\newcommand{\ie}{i.e.\ }
\newcommand{\cf}{cf.\ }
\newcommand{\resp}{respectively\ }
\newcommand{\Subsection}{\subsection}
\providecommand{\nobreakdash}{\nobreak\hbox{-}}
\setlength{\emergencystretch}{2em}
\multlinegap0pt
\usepackage{amssymb}
\usepackage{mathtools}
\usepackage{float}
\usepackage{tikz}
\newtheorem{theorem}{Theorem}
\title{Minimal spheres and scalar curvature}
\author{Talant Talipov}
\date{Source: arXiv:2605.21607v1 (CC BY 4.0)}
\begin{document}
\maketitle

\noindent\emph{Source and licence.} Minimal spheres and scalar curvature. Version arXiv:2605.21607v1 (CC BY 4.0), distributed by arXiv under the Creative Commons Attribution 4.0 International licence, http://creativecommons.org/licenses/by/4.0/, as recorded in the arXiv licence metadata for this exact version and shown on the abstract page https://arxiv.org/abs/2605.21607v1. Full licence notice: \texttt{licenses/arxiv-2605.21607-CC-BY-4.0.txt}.\par\smallskip

\noindent\textit{Editorial scope.} Counted unit: the article's theorem thm:yau (source0.tex lines 99--102). The article's complete original proof of it is delivered with it (source0.tex lines 673--791, one \texttt{proof} environment of 119 lines, which the article places away from the statement). No mathematics is written, repaired or re-derived: the statement and the proof are the article's own lines, in the article's own notation, and no retained line is substituted or re-flowed.  Retained uncounted as the source-local context the proof needs: lines 67--77.  Nothing was omitted from any retained span, so the package carries no omission record.  No retained line contains a citation, so the wrapper carries no bibliography and no bibliography entry.  No retained line contains a figure environment, an image-inclusion command or drawing code, so no image is delivered and the package declares no asset dependency.  Editorial formatting only, in addition: an AMS \texttt{amsart} preamble that restates the article's own declarations and macros used by the retained lines, namely the environments \texttt{theorem} on separate counters, together with the standard packages those lines need.  The retained environments print in this document as Theorem 1 in order of appearance, whereas the article prints the same items with the numbering of its own full document.  The complete native source of the article as distributed in the arXiv e-print for this exact version is bundled unchanged as \texttt{sources/201071/source0.tex}.
\begin{theorem}
\label{thm:main}
    Let \(\Lambda_0>0\). Suppose that \((S^3,g_0)\) satisfies \(\operatorname{Ric}_{g_0}>0\) and \(R_{g_0}\ge \Lambda_0\). Then there exist four distinct embedded minimal two-spheres
    \[
    \Sigma_1,\ldots,\Sigma_4 \subset (S^3,g_0)
    \]
    such that, for every \(i=1,\ldots,4\),
    \[
    \operatorname{area}_{g_0}(\Sigma_i)\le \frac{12\pi(i+1)}{\Lambda_0}.
    \]
\end{theorem}

\begin{theorem}
\label{thm:yau}
For fixed $b,c$ and $d$, if $a$ is chosen sufficiently large, then the ellipsoid $E(a,b,c,d)$ contains at least three non-planar embedded minimal two-spheres.
\end{theorem}

\begin{proof}[Proof of Theorem~\ref{thm:yau}]
Fix \(b>c>d>0\), and write
\[
E_a:=E(a,b,c,d).
\]
Let \(g_a\) denote the metric induced on \(E_a\) by the Euclidean metric of
\(\mathbb R^4\). For \(j=1,\ldots,4\), denote the coordinate planar minimal
two-spheres by
\[
P_j(a):=E_a\cap\{x_j=0\}.
\]

We first record a scalar curvature lower bound for \(E_a\), uniform for all
\(a>b\). Let
\[
A_a=\operatorname{diag}(a^{-2},b^{-2},c^{-2},d^{-2}).
\]
Then
\[
E_a=\{x\in\mathbb R^4:\langle A_a x,x\rangle=1\}.
\]
For \(v\in T_xE_a\), the second fundamental form is
\[
\mathrm{II}(v,v)=\frac{\langle A_a v,v\rangle}{|A_a x|}.
\]
Moreover,
\[
|A_a x|^2
=
\langle A_a^2x,x\rangle
\le
d^{-2}\langle A_a x,x\rangle
=
d^{-2}.
\]
Hence,
\[
|A_a x|\le d^{-1}.
\]

Let \(H_1=\{v_1=0\}\subset\mathbb R^4\). Since
\(\dim T_xE_a=3\) and \(\dim H_1=3\), we have
\[
\dim(T_xE_a\cap H_1)\ge2.
\]
For every \(v\in T_xE_a\cap H_1\),
\[
\langle A_a v,v\rangle
\ge
b^{-2}|v|^2.
\]
Then
\[
\mathrm{II}(v,v)
\ge
\frac{b^{-2}}{|A_a x|}|v|^2
\ge
\frac{d}{b^2}|v|^2.
\]
By the min-max characterization of eigenvalues, at least two principal curvatures
of \(E_a\) at \(x\) are bounded below by \(d/b^2\). Since \(E_a\) is strictly convex,
all principal curvatures are positive. Thus the Gauss equation gives
\[
R_{g_a}
=
2\sum_{1\le \alpha<\beta\le3}\kappa_\alpha\kappa_\beta
\ge
\frac{2d^2}{b^4} =: \Lambda_0.
\]
Also, since all principal curvatures of \(E_a\) are positive, \(g_a\) has positive
Ricci curvature.

Applying Theorem~\ref{thm:main} to \((E_a,g_a)\), we obtain four distinct embedded minimal two-spheres
\[
\Sigma_1(a),\ldots,\Sigma_4(a)\subset E_a
\]
such that
\[
\operatorname{area}_{g_a}(\Sigma_i(a))
\le
\frac{12\pi(i+1)}{\Lambda_0}
\le
\frac{60\pi}{\Lambda_0}
=
\frac{30\pi b^4}{d^2}
\]
for every \(i=1,\ldots,4\).

We now compare this uniform area bound with the areas of the planar spheres. The planar sphere \(P_1(a)\) is independent of \(a\), and hence has bounded area. On the other hand, for \(j=2,3,4\), the planar sphere \(P_j(a)\) is a two-dimensional ellipsoid with one semi-axis equal to \(a\). Consider the restriction of the coordinate function \(x_1\) to \(P_j(a)\).
Since \(|\nabla^{P_j(a)}x_1|\le1\), by the coarea formula,
\[
\operatorname{area}(P_j(a))
\ge
\int_{-a/2}^{a/2}
\mathcal H^1\bigl(P_j(a)\cap\{x_1=t\}\bigr)\,dt.
\]
For \(|t|\le a/2\), the curve \(P_j(a)\cap\{x_1=t\}\) is an ellipse with smaller semi-axis at least
\[
d\sqrt{1-\frac{t^2}{a^2}}
\ge
\frac{\sqrt3}{2}d.
\]
Therefore its length is at least \(\sqrt3\pi d\), and then
\[
\operatorname{area}(P_j(a))
\ge
\sqrt3\pi da
\]
for \(j=2,3,4\). Thus, for \(a\) sufficiently large,
\[
\operatorname{area}(P_j(a))
>
\frac{30\pi b^4}{d^2}
\]
for \(j=2,3,4\).

For such \(a\), none of the four minimal two-spheres
\(\Sigma_1(a),\ldots,\Sigma_4(a)\) can be equal to \(P_2(a),P_3(a)\), or \(P_4(a)\), since each \(\Sigma_i(a)\) has area at most \(30\pi b^4/d^2\), while each of \(P_2(a),P_3(a),P_4(a)\) has strictly larger area. Since the four spheres \(\Sigma_1(a),\ldots,\Sigma_4(a)\) are distinct, at most one of them can coincide with the remaining planar sphere \(P_1(a)\). Hence at least three of the four minimal spheres \(\Sigma_1(a),\ldots,\Sigma_4(a)\) are non-planar.
\end{proof}

\end{document}
